\section*{Capítulo \ref{ch:g1:materials-around-us} --- ¿De qué están hechas las cosas?}

\begin{solution}{exo:g1:materials-around-us:1}
Una silla es un \omterm{def:g1:materials-around-us:object}{objeto}. La madera es un \omterm{def:g1:materials-around-us:material}{material}: con ella se puede
hacer una silla.
\end{solution}

\begin{solution}{exo:g1:materials-around-us:2}
El cristal de la ventana: vidrio. El clavo: metal. El guijarro: piedra.
El libro: papel. El calcetín: tela (lana o algodón).
\end{solution}

\begin{solution}{exo:g1:materials-around-us:3}
De madera: un lápiz, una puerta (o una silla, una mesa). De plástico:
un cuenco, un cubo (o un juguete, una botella).
\end{solution}

\begin{solution}{exo:g1:materials-around-us:4}
El tenedor de metal: es el único que no es ni de madera ni de plástico.
(Dos son de madera y uno es de plástico.)
\end{solution}

\begin{solution}{exo:g1:materials-around-us:5}
El vidrio: es transparente.
\end{solution}

\begin{solution}{exo:g1:materials-around-us:6}
El clavo de acero. Un imán atrae el hierro y el acero, no el vidrio ni
la madera.
\end{solution}

\begin{solution}{exo:g1:materials-around-us:7}
Flotan dos \omterm{def:g1:materials-around-us:object}{objetos} (el taco de madera y el tapón de plástico); se
hunden tres (el guijarro, la canica y el clavo).
\end{solution}

\begin{solution}{exo:g1:materials-around-us:8}
Las hojas son de metal y los mangos, de plástico. Las tijeras están
hechas de dos \omterm{def:g1:materials-around-us:material}{materiales}.
\end{solution}

\begin{solution}{exo:g1:materials-around-us:9}
En el aro azul: es de metal. No está en el aro naranja porque un imán
no la atrae: es de aluminio, no de hierro ni de acero.
\end{solution}

\begin{solution}{exo:g1:materials-around-us:10}
Una ventana tiene que dejar entrar la luz, y tenemos que ver a través
de ella. El vidrio es transparente; la madera, no.
\end{solution}

\begin{solution}{exo:g1:materials-around-us:11}
``Sí'': el clavo, el clip y la llave, 3 \omterm{def:g1:materials-around-us:object}{objetos}. ``No'': el guijarro,
la cuchara, la lata, la canica y el calcetín, 5 \omterm{def:g1:materials-around-us:object}{objetos}.
$3 + 5 = 8$: cada \omterm{def:g1:materials-around-us:object}{objeto} tiene su sitio.
\end{solution}
